%*******************************************************************************
%*********************************** First Chapter *****************************
%*******************************************************************************
\graphicspath{{Figures}}
\epigraphhead[550]{"This is an awesome quote"
\par\hfill\textsc{― Random Person}}


\part{Introduction}  %Title of the First Chapter

% \cleardoublepage
\chapter{Introduction}
\label{ch:intro}



\section{Context}
\label{sec:context}

Since the Industrial Revolution, technological innovations have profoundly transformed production, transport, and everyday life. Mechanization and the rise of mass manufacturing improved living standards and enabled unprecedented population growth. Moreover, the development of motor vehicles, railways, and later air transport allowed goods and people to travel farther and faster than ever before. Yet, this progress--whose pace shows little sign of slowing--has also brought significant unintended consequences~\cite{IPCC2023, UNEP2023}, leading to patterns of resource consumption, pollution, and land use that increasingly challenge the sustainability of our society. Examples include large-scale pollutant emissions, deforestation, and growing volumes of waste, all of which place great pressure on the environment and human health.

These patterns have become particularly visible in cities, where rapid population growth has intensified the demand for housing, services, and infrastructure. This expansion has created enormous challenges in managing mobility and urban development. Even though cities occupy just 3\% of the Earth's land, they account for 60–80\% of energy consumption and 75\% of carbon emissions~\cite{UN_SDG11_2015}. The \acf{sdg} recognize the urgency of this situation: Goal 11 aims to make cities inclusive, safe, resilient, and sustainable, emphasizing the need to transform both urban spaces and mobility systems. Yet, in 2022, only half of the urban population had convenient access to public transport, and this gap is expected to widen further: by 2050, around 70\% of the world’s population is projected to live in urban areas, compared with about 57\% today~\cite{UN_SDG11_2015}. Together, these factors highlight why rethinking how cities are planned and managed is essential.

Looking more specifically at urban mobility, the rapid motorization of cities has created profound challenges for planners and policymakers. In many metropolitan areas, private car use, while often convenient and flexible for individuals, has been systematically prioritized over sustainable alternatives. Decades of urban development have favored expanding road capacity, allocating extensive land to parking, and promoting low-density suburban growth instead of compact, mixed-use neighborhoods. These choices have resulted in transport systems that generate a disproportionate share of greenhouse gas emissions and local air pollution, contribute to chronic congestion, and undermine the safety and quality of urban public space. In dense urban environments, where exposure to vehicle emissions and noise is highest, these impacts have direct consequences for public health and social equity. In response, over the last decade, many cities have launched ambitious policies and investments to reverse these trends, including better public transport, cycling infrastructure, greener public spaces, measures to curb car use, and expanded shared mobility services~\cite{UNHabitat2020}. Despite progress, the transition toward sustainable urban mobility remains incomplete.

Among these initiatives, \acf{bss}s have emerged as one of the most visible examples of this shift. By providing flexible, low-carbon transport options, they have contributed to reducing car dependency and promoting active mobility. In many cities, \acf{bss}s are integrated with public transport networks and can help improve first- and last-mile connectivity, ease congestion, and support healthier lifestyles~\cite{Fishman2016, Shaheen2010}. However, their success depends on careful planning and effective operations. Challenges such as station siting, fleet rebalancing, vehicle maintenance, and ensuring equitable access require well-designed management approaches.

\acf{ai} and advanced data analytics have increasingly been recognized as powerful tools for addressing the complexity of modern urban mobility systems. The proliferation of sensors, mobile devices, and digital platforms has generated unprecedented volumes of mobility data, enabling more accurate modeling, prediction, and optimization of transport services. In the context of \acf{bss}s, \acf{ai} methods such as demand forecasting, predictive maintenance, dynamic rebalancing, and user behavior analysis can help operators improve efficiency, reduce operational costs, and enhance user satisfaction~\cite{Zhang2021, Luo2022}. More broadly, \acf{ml} and optimization techniques are being applied to a wide range of transport challenges, including real-time traffic management, multimodal trip planning, and the design of infrastructure that supports more sustainable and equitable mobility patterns. However, the integration of these methods into operational decision-making requires careful consideration of data quality, computational constraints, interpretability, and alignment with policy goals. Understanding how these tools can enhance \acf{bss} operations is essential for designing sustainable, resilient urban mobility systems.


This thesis explores how \acf{ai} and data-driven approaches can enhance the planning and operation of \acf{bss}s as part of broader sustainable mobility strategies. Focusing on the case of Barcelona and its \acf{bss} system Bicing, it develops and evaluates predictive models and optimization methods to address key challenges such as demand forecasting, battery level estimation, and the strategic placement of stations. Through these use cases, the research aims to demonstrate how \ac{ai} can support more efficient, equitable, and resilient \ac{bss} services, ultimately contributing to healthier and more sustainable urban environments.




















